\subsection{Objective}

Task 2 requires implementing a $5\times5$ averaging filter, applying it to the
image, trying additional filter sizes, and comparing the smoothing results.
Because Task 1 selected the 2--98\% percentile-stretched image as the preferred
input for the downstream stages, every smoothing experiment here operates on
that enhanced grayscale image.

\subsection{Implemented Averaging Filter}

An averaging (box) filter~\cite{gonzalez2018} of size $s\times s$ replaces each pixel with the mean
of the $s\times s$ neighbourhood centred on it:
\begin{equation}
  I_{\mathrm{out}}(x,y) =
  \frac{1}{s^{2}}
  \sum_{i=-r}^{r}\sum_{j=-r}^{r} I_{\mathrm{in}}(x+i,y+j),
  \qquad r=\frac{s-1}{2}.
  \label{eq:task2-box}
\end{equation}
The implementation exploits the separability of the box kernel and computes the
sliding sums with a summed-area table, so the running time is $O(hw)$
independent of $s$ rather than $O(hw s^{2})$ for a naive sliding window.  The
filter is implemented from first principles with NumPy only in
\texttt{common.box\_filter}; boundary pixels are handled by replicating the
edge values.  The required $5\times5$ size was implemented together with four
additional odd sizes, $3\times3$, $7\times7$, $11\times11$, and $15\times15$, to
make the smoothing--detail trade-off visible.

\begin{figure}[!htbp]
  \centering
  \includegraphics[width=\textwidth]{task2_comparison.png}
  \caption{The enhanced image and the box-filtered results for five filter
  sizes.  All images use the same display range of 0--255.}
  \label{fig:task2-comparison}
\end{figure}

\begin{figure}[!htbp]
  \centering
  \includegraphics[width=0.9\textwidth]{task2_zoom.png}
  \caption{Zoomed view of the \texttt{HD4} characters.  Larger kernels remove
  more background texture but progressively blur the stroke edges.}
  \label{fig:task2-zoom}
\end{figure}

\subsection{Evaluation Protocol}

The smoothing results were compared both visually and quantitatively.  Three
values were measured on fixed image regions so that all filter sizes are judged
identically:

\begin{itemize}
  \item \textbf{Background standard deviation}: the pooled standard deviation
  of five fixed character-free background patches.  A smaller value means the
  ring-like texture and speckle noise were suppressed more strongly.
  \item \textbf{Foreground gradient}: the mean gradient magnitude over five
  fixed character-stroke patches.  A smaller value means the stroke edges were
  blurred more, i.e.\ more detail was lost.
  \item \textbf{PSNR}: peak signal-to-noise ratio relative to the unfiltered
  enhanced image, measuring how much the whole image changed.
\end{itemize}

\begin{table}[!htbp]
  \centering
  \caption{Task 2 quantitative comparison of the averaging-filter sizes.}
  \label{tab:task2-metrics}
  \begin{tabular}{lccc}
    \toprule
    Filter & Background std. & Foreground gradient & PSNR (dB) \\
    \midrule
    None (enhanced baseline) & 33.42 & 20.18 & $\infty$ \\
    $3\times3$ box & 29.56 & 14.75 & 28.50 \\
    $5\times5$ box & 26.46 & 11.00 & 24.03 \\
    $7\times7$ box & 24.09 & 9.44 & 21.95 \\
    $11\times11$ box & 19.91 & 7.86 & 20.17 \\
    $15\times15$ box & 16.36 & 6.75 & 18.90 \\
    \bottomrule
  \end{tabular}
\end{table}

\subsection{Result Analysis}

The measurements in \cref{tab:task2-metrics} show a clear monotonic trade-off.
As the kernel grows, the background standard deviation falls from 33.42 to
16.36, confirming stronger noise suppression, but the foreground gradient also
falls from 20.18 to 6.75, meaning the character edges are increasingly blurred.
The $3\times3$ filter is too weak: it reduces the background noise by only
11.5\% while the largest $15\times15$ filter removes 51.0\% of the noise but
erases 66.5\% of the edge sharpness.  In \cref{fig:task2-zoom}, the thin
strokes of the characters visibly soften at $11\times11$ and start to blend
with the background at $15\times15$; such damage would later break or merge
character strokes during thresholding.

The required $5\times5$ filter is the best compromise for this image.  It
reduces the background standard deviation by 20.8\% while retaining 54.5\% of
the edge gradient magnitude, so isolated speckles are suppressed without
destroying the thin character strokes.  Its PSNR of 24.03\,dB confirms that the
change is moderate.  This $5\times5$ smoothing was therefore reused in
Task~3 for row localization and in Task~4 before thresholding, keeping the
image-processing pipeline consistent.

\FloatBarrier
